%% ============================================================
%% Chapter 3 — Methodology
%% ============================================================
\chapter{Methodology}
\label{chap:methodology}

\section{Dataset Preparation}
\label{sec:meth_dataset}

\subsection{Source Data}

We work with the official SQuAD~2.0 JSON distribution
\cite{rajpurkar2018squad2}. The training file contains 130,319 question–answer
pairs spanning 442 Wikipedia articles; the development file contains 11,873
pairs. Each record specifies the article title, the passage (paragraph) text,
the question, and either a list of answer spans (possibly empty for unanswerable
questions) alongside a boolean \texttt{is\_impossible} flag.

\subsection{Article-Level Splitting}
\label{sec:meth_split}

A subtle but critical experimental design decision concerns how to partition the
data into training, validation, and test sets. A naive random split at the
question level risks placing questions about the same article—and even the same
passage—in both training and test sets. Because BM25 builds an inverted index
over the passages, a passage seen at training time will produce identical BM25
features at test time, inflating test-set performance through data leakage.

To prevent this, we perform an \emph{article-level} split: all questions
belonging to a given article are assigned exclusively to one split. We
stratify on article-level answerability to maintain an equal proportion of
answerable and unanswerable questions in each split. The final partition
contains 500 questions distributed as follows (see
Table~\ref{tab:dataset_stats} and Figure~\ref{fig:dataset_distribution}):

\begin{itemize}
  \item \textbf{Train:} 300 questions — 150 answerable, 150 unanswerable
        (50\% / 50\%, perfectly balanced; 176 unique articles).
  \item \textbf{Validation:} 100 questions — 50 answerable, 50 unanswerable
        (50\% / 50\%, perfectly balanced; 61 unique articles).
  \item \textbf{Test:} 100 questions — 50 answerable, 50 unanswerable
        (50\% / 50\%, perfectly balanced; 55 unique articles).
\end{itemize}

The test set is kept entirely unseen during hyper-parameter tuning; all
threshold and classifier decisions are made based solely on validation-set
performance.

\section{BM25 Retrieval and Feature Extraction}
\label{sec:meth_bm25}

\subsection{Indexing}

For each data split, we build a BM25 index over the passages belonging to that
split. Using default Okapi BM25 parameters ($k_1 = 1.5$, $b = 0.75$), we
tokenise passage text with whitespace tokenisation after lower-casing. Each
passage is treated as an independent document in the index; articles may contain
multiple passages.

\subsection{Top-$k$ Retrieval}

Given a question $q$, we query the split's BM25 index and retrieve the top-$k$
passages ranked by BM25 score. We use $k = 5$ throughout. The top-ranked
passage is passed to the QA reader; the remaining passages are used only for
feature computation.

\subsection{Retrieval Features}

From the BM25 retrieval results for each question, we extract four inference-ready
features:

\begin{description}
  \item[\texttt{top\_bm25\_score}] The BM25 score of the first-ranked passage.
        This is the primary indicator of retrieval confidence.

  \item[\texttt{second\_bm25\_score}] The BM25 score of the second-ranked
        passage. The relative magnitudes of ranks 1 and 2 convey whether the
        top result is unambiguously best or faces strong competition.

  \item[\texttt{bm25\_score\_gap}] The difference
        $\texttt{top\_bm25\_score} - \texttt{second\_bm25\_score}$. A large
        gap indicates that the top passage is a qualitatively better match than
        all alternatives.

  \item[\texttt{mean\_topk\_bm25}] The arithmetic mean of BM25 scores across
        the top-$k = 5$ retrieved passages. A high mean suggests that the query
        matches many passages, which may indicate a broad topic or retrieval
        ambiguity.
\end{description}

\paragraph{Offline Diagnostic Feature (\texttt{source\_passage\_retrieved}).}
In addition to the operational features above, we compute a binary indicator
whether the original SQuAD source passage appears in the top-$k$ retrieved IDs.
Critically, because the true source passage ID is an artifact of the benchmark
annotation and is completely unavailable for new user questions at test time,
this feature is \emph{strictly excluded} from the classifier's predictive
feature set to eliminate data leakage and guarantee training/inference parity.
It is retained solely for offline retrieval diagnostics.

\section{QA Reader}
\label{sec:meth_reader}

\subsection{Model}

We use the \texttt{deepset/roberta-base-squad2} checkpoint from the
Hugging Face model hub, a RoBERTa-base \cite{liu2019roberta} model fine-tuned
on SQuAD~2.0. The model processes a concatenated \texttt{[CLS] question [SEP]
passage [SEP]} sequence and outputs a distribution over start and end positions,
together with a CLS-based no-answer score.

\subsection{Inference Protocol}

We set \texttt{handle\_impossible\_answer=True}, instructing the pipeline to
compute both a span score and a no-answer probability. For each question, we
pass only the \emph{top-1} BM25-retrieved passage to the reader (to simulate a
standard retrieval–reader pipeline). The reader produces:

\begin{description}
  \item[\texttt{qa\_score}] The probability assigned to the extracted answer
        span. High values indicate reader confidence.

  \item[\texttt{qa\_no\_answer\_probability}] The probability that no answer
        exists in the passage. This is derived from the CLS token logit and is
        directly trained by the SQuAD~2.0 fine-tuning objective.
\end{description}

\subsection{Reader Feature Extraction}

In addition to the two QA probabilities, we derive:

\begin{description}
  \item[\texttt{predicted\_answer\_length}] The character length of the
        extracted answer span (0 for empty answers).

  \item[\texttt{answer\_nonempty}] A binary flag: 1 if the reader returned a
        non-empty answer string; 0 if it abstained at the extraction level.

  \item[\texttt{question\_length}] The character length of the question string.

  \item[\texttt{question\_token\_count}] The number of whitespace-delimited
        tokens in the question.
\end{description}

\section{Abstention Strategies}
\label{sec:meth_strategies}

We now describe the three abstention strategies in detail.

\subsection{Strategy 1 — Always Answer (Baseline)}
\label{sec:meth_always_answer}

The baseline strategy applies \emph{no abstention}: the system always returns
the reader's top-1 extracted span, regardless of retrieval or reader scores.
For questions where the reader itself produces an empty answer (triggered by
\texttt{handle\_impossible\_answer=True}), the system returns an empty string.
This strategy achieves maximum coverage (all questions receive a response) at
the cost of precision, as it cannot suppress false answers.

Formally, the selection function is $g_\text{base}(x) = 1$ for all inputs $x$.

\subsection{Strategy 2 — BM25 Threshold}
\label{sec:meth_threshold}

The BM25 Threshold strategy abstains on questions whose top BM25 score falls
below a threshold $\tau$:

\begin{equation}
  g_\tau(x) =
  \begin{cases}
    1 & \text{if } \texttt{top\_bm25\_score}(x) \geq \tau, \\
    0 & \text{otherwise (abstain).}
  \end{cases}
  \label{eq:bm25_threshold}
\end{equation}

The threshold $\tau$ is selected to maximise $\text{F1}$ on the validation set
over the grid $\tau \in \{5, 6, \ldots, 50\}$ (integer values). The selected
threshold is then applied \emph{unchanged} to the test set, ensuring no
information from the test set influences the threshold selection.

The intuition is that a low BM25 score indicates poor lexical overlap between
the question and all retrieved passages, suggesting the corpus does not contain
relevant information and the system should abstain.

\subsection{Strategy 3 — Answerability Classifier}
\label{sec:meth_classifier}

The Answerability Classifier is a logistic regression model trained on the
ten inference-ready features enumerated in Table~\ref{tab:features}. It outputs a probability
$\hat{p}(x)$ that the question is answerable given the retrieved context. The
selection function is:

\begin{equation}
  g_\theta(x) =
  \begin{cases}
    1 & \text{if } \hat{p}(x) \geq \delta, \\
    0 & \text{otherwise (abstain),}
  \end{cases}
  \label{eq:classifier_threshold}
\end{equation}

where $\delta$ is a probability threshold tuned on the validation set by
maximising validation F1. The optimal threshold found is $\delta^* = 0.42$.

\subsubsection{Pre-processing Pipeline}

Training instances are drawn from the training split. All numerical features
are standardised with \texttt{sklearn}'s \texttt{StandardScaler}; missing values
(which arise when fewer than $k$ passages are retrieved) are imputed with
\texttt{SimpleImputer} using the median strategy. The pipeline is composed as:
\texttt{SimpleImputer} $\rightarrow$ \texttt{StandardScaler} $\rightarrow$ \texttt{LogisticRegression(class\_weight='balanced')}.

\subsubsection{Feature Set}

\begin{table}[htbp]
  \centering
  \caption{Inference feature set for the Answerability Classifier (10 features).}
  \label{tab:features}
  \begin{tabular}{lll}
    \toprule
    \textbf{Feature} & \textbf{Source} & \textbf{Description} \\
    \midrule
    \texttt{top\_bm25\_score}        & BM25 & Score of rank-1 retrieved passage \\
    \texttt{second\_bm25\_score}     & BM25 & Score of rank-2 retrieved passage \\
    \texttt{bm25\_score\_gap}        & BM25 & Rank-1 minus rank-2 BM25 scores \\
    \texttt{mean\_topk\_bm25}        & BM25 & Mean score over top-$k$ passages \\
    \texttt{qa\_score}               & Reader & Reader confidence in answer span \\
    \texttt{qa\_no\_answer\_probability} & Reader & Reader no-answer probability \\
    \texttt{predicted\_answer\_length} & Reader & Character length of predicted span \\
    \texttt{answer\_nonempty}        & Reader & 1 if reader returned non-empty answer \\
    \texttt{question\_length}        & Query & Character length of question \\
    \texttt{question\_token\_count}  & Query & Token count of question \\
    \bottomrule
  \end{tabular}
\end{table}

\subsubsection{Training and Threshold Tuning}

Logistic regression is trained with balanced class weights and L2 regularisation.
After training on the 300 training questions, we sweep $\delta$ over $[0.01, 0.99]$ in steps
of $0.01$ on the 100 validation questions, computing precision, recall, and F1. The threshold
that maximises validation F1 is retained as $\delta^* = 0.42$ (Validation F1 = 0.6912) and applied to the test set.

\subsection{Evaluation Metrics}
\label{sec:meth_metrics}

Let $A$ be the set of questions for which the system answers ($g(x)=1$), $C$ be
the set of questions where the system's answer is \emph{correct} (exact match
or F1 above a threshold), $U$ be the truly unanswerable questions, and $N$
be the total number of questions. We report:

\begin{align}
  \text{Precision} &= \frac{|C|}{|A|},
  \label{eq:precision} \\[4pt]
  \text{Recall} &= \frac{|C \cap \text{Answerable}|}{|\text{Answerable}|},
  \label{eq:recall} \\[4pt]
  \text{F1} &= 2 \cdot \frac{\text{Precision} \cdot \text{Recall}}
                              {\text{Precision} + \text{Recall}},
  \label{eq:f1} \\[4pt]
  \text{Coverage} &= \frac{|A|}{N},
  \label{eq:coverage} \\[4pt]
  \text{False-Answer Rate} &= \frac{|\{x \in A : x \in U\}|}{|U|},
  \label{eq:far} \\[4pt]
  \text{False-Abstention Rate} &= \frac{|\{x \notin A : x \in \text{Answerable}\}|}{|\text{Answerable}|}.
  \label{eq:fab}
\end{align}

Precision and F1 reward the system for giving correct answers when it does
answer. Coverage measures how much of the question set the system addresses.
False-Answer Rate and False-Abstention Rate characterise the two types of
errors that result from imperfect abstention.
